    \chapter{Mathematical Prizes and Congresses}

    Science Quiz often asks for the short association between a prize, a
    congress, a laureate, and a mathematical achievement. This chapter records
    the facts that should be recognized quickly: not full biographies, but the
    names, years, keywords, and historical landmarks most likely to become quiz
    answers.

    \begin{fact}[Science Quiz Prize Triad]
        The three mathematics prizes that should be recognized first are
        \[
            \text{Fields Medal},\qquad
            \text{Abel Prize},\qquad
            \text{Wolf Prize in Mathematics}.
        \]
        Memorize each prize through the association
        \[
            \text{prize}
            \quad\longleftrightarrow\quad
            \text{laureate}
            \quad\longleftrightarrow\quad
            \text{achievement or area}.
        \]
    \end{fact}

    \begin{strategy}
        History questions often ask for a short association rather than a full
        biography. Current prominence, a recent award, or a connection to the
        \POSTECH--\KAIST{} context makes such an association especially relevant.
    \end{strategy}

    \section*{Fields Medal and ICM}

    The International Congress of Mathematicians, commonly abbreviated as ICM,
    is the most important international congress in mathematics. The Fields
    Medal is awarded by the International Mathematical Union at the ICM every
    four years. It is named after John Charles Fields. A Fields Medalist must
    satisfy the age condition: the candidate's $40$th birthday must not occur
    before January $1$ of the year of the Congress.

    \vspace{1em}

    {\small
    \renewcommand{\arraystretch}{1.18}
    \begin{longtable}{@{}L{0.07\textwidth}L{0.16\textwidth}L{0.33\textwidth}L{0.36\textwidth}@{}}
        \toprule
        Year & ICM location & Medalists & Quiz keywords \\
        \midrule
        \endfirsthead

        \toprule
        Year & ICM location & Medalists & Quiz keywords \\
        \midrule
        \endhead

        \midrule
        \multicolumn{4}{r}{Continued on the next page.} \\
        \endfoot

        \bottomrule
        \endlastfoot

        1936
        & Oslo, Norway
        & Lars Ahlfors; Jesse Douglas
        & Complex analysis, Riemann surfaces; minimal surfaces, Plateau problem. \\

        1950
        & Cambridge, USA
        & Laurent Schwartz; Atle Selberg
        & Distribution theory; analytic number theory, sieve methods, zeta functions. \\

        1954
        & Amsterdam, Netherlands
        & Kunihiko Kodaira; Jean-Pierre Serre
        & Complex geometry, algebraic geometry; algebraic topology, sheaf theory. \\

        1958
        & Edinburgh, UK
        & Klaus Roth; René Thom
        & Diophantine approximation; cobordism, differential topology. \\

        1962
        & Stockholm, Sweden
        & Lars Hörmander; John Milnor
        & Partial differential equations; differential topology, exotic spheres. \\

        1966
        & Moscow, USSR
        & Michael Atiyah; Paul Cohen; Alexander Grothendieck; Stephen Smale
        & K-theory and index theorem; forcing and continuum hypothesis;
          modern algebraic geometry; generalized Poincaré conjecture. \\

        1970
        & Nice, France
        & Alan Baker; Heisuke Hironaka; Serge Novikov; John G. Thompson
        & Transcendental number theory; resolution of singularities;
          topology and geometry; finite group theory. \\

        1974
        & Vancouver, Canada
        & Enrico Bombieri; David Mumford
        & Number theory, complex analysis, minimal surfaces; moduli spaces and
          algebraic geometry. \\

        1978
        & Helsinki, Finland
        & Pierre Deligne; Charles Fefferman; Grigory Margulis; Daniel Quillen
        & Weil conjectures; several complex variables; Lie groups and ergodic
          theory; algebraic K-theory. \\

        1982
        & Warsaw, Poland (held in 1983)
        & Alain Connes; William Thurston; Shing-Tung Yau
        & Noncommutative geometry; $3$-manifolds and low-dimensional topology;
          Calabi conjecture and geometric analysis. \\

        1986
        & Berkeley, USA
        & Simon Donaldson; Gerd Faltings; Michael Freedman
        & Gauge theory and $4$-manifolds; Mordell conjecture and arithmetic
          geometry; topological $4$-dimensional Poincaré conjecture. \\

        1990
        & Kyoto, Japan
        & Vladimir Drinfeld; Vaughan Jones; Shigefumi Mori; Edward Witten
        & Quantum groups and Langlands program; Jones polynomial; minimal
          model program; mathematical physics and quantum field theory. \\

        1994
        & Zurich, Switzerland
        & Jean Bourgain; Pierre-Louis Lions; Jean-Christophe Yoccoz; Efim Zelmanov
        & Analysis and Banach spaces; partial differential equations; dynamical
          systems; restricted Burnside problem. \\

        1998
        & Berlin, Germany
        & Richard Borcherds; W. Timothy Gowers; Maxim Kontsevich; Curtis T. McMullen
        & Moonshine and automorphic forms; functional analysis and combinatorics;
          algebraic geometry and mathematical physics; complex dynamics. \\

        2002
        & Beijing, China
        & Laurent Lafforgue; Vladimir Voevodsky
        & Langlands correspondence over function fields; motivic cohomology and
          $\mathbb{A}^1$-homotopy theory. \\

        2006
        & Madrid, Spain
        & Andrei Okounkov; Grigori Perelman; Terence Tao; Wendelin Werner
        & Probability, representation theory, algebraic geometry; Ricci flow;
          PDE, harmonic analysis, additive combinatorics; SLE and Brownian motion. \\

        2010
        & Hyderabad, India
        & Elon Lindenstrauss; Ngo Bao Chau; Stanislav Smirnov; Cédric Villani
        & Ergodic theory and number theory; Fundamental Lemma; conformal
          invariance; Boltzmann equation and optimal transport. \\

        2014
        & Seoul, South Korea
        & Artur Avila; Manjul Bhargava; Martin Hairer; Maryam Mirzakhani
        & Dynamical systems; geometry of numbers; stochastic PDE;
          Riemann surfaces and moduli spaces. \\

        2018
        & Rio de Janeiro, Brazil
        & Caucher Birkar; Alessio Figalli; Peter Scholze; Akshay Venkatesh
        & Fano varieties and minimal model program; optimal transport;
          $p$-adic geometry; number theory and homogeneous dynamics. \\

        2022
        & Virtual; awards in Helsinki, Finland
        & Hugo Duminil-Copin; June Huh; James Maynard; Maryna Viazovska
        & Phase transitions in statistical physics; Hodge theory in combinatorics
          and matroids; prime numbers; sphere packing and the $E_8$ lattice. \\

        2026
        & Philadelphia, USA
        & Yu Deng; John Pardon; Jacob Tsimerman; Hong Wang
        & PDE and kinetic theory; symplectic geometry and topology;
          arithmetic geometry and o-minimality; harmonic analysis, geometric
          measure theory, and the Kakeya problem. \\
    \end{longtable}
    }

    \begin{fact}[ICM and Fields Medal landmarks]
        The following associations are especially useful.
        \begin{itemize}
            \item ICM stands for International Congress of Mathematicians.
            \item IMU stands for International Mathematical Union.
            \item The Fields Medal was first awarded in 1936 at ICM Oslo.
            \item The first two Fields Medalists were Lars Ahlfors and Jesse Douglas.
            \item The first year with four Fields Medalists was 1966.
            \item Jean-Pierre Serre is the youngest Fields Medalist.
            \item Edward Witten was the first physicist to receive the Fields Medal; as an
            undergraduate, he majored in history rather than mathematics.
            \item Andrew Wiles received a special IMU silver plaque in 1998 for
                  his proof of Fermat's Last Theorem.
            \item Grigori Perelman declined the Fields Medal in 2006.
            \item Maryam Mirzakhani was the first woman Fields Medalist.
            \item Maryna Viazovska was the second woman Fields Medalist.
            \item ICM 2014 was held in Seoul, South Korea.
            \item ICM 2022 was held as a virtual congress; the prize ceremony was
                  held in Helsinki, Finland.
            \item ICM 2026 was held in Philadelphia, USA.
            \item ICM 2030 is scheduled to be held in Glasgow, UK.
        \end{itemize}
    \end{fact}

    \begin{fact}[Fields Medal design]
        The obverse of the Fields Medal bears a profile of Archimedes and the
        Latin inscription
        \[
            \text{\textit{TRANSIRE SUUM PECTUS MUNDOQUE POTIRI}}.
        \]
        The reverse shows a sphere inscribed in a cylinder, recalling the result
        that Archimedes requested for his tomb. The cylinder and the inscribed
        sphere have volume ratio
        \[
            3:2.
        \]
    \end{fact}

    \section*{Abel Prize}

    The Abel Prize is awarded by the Norwegian Academy of Science and Letters.
    It is named after Niels Henrik Abel and was first awarded in 2003. Unlike
    the Fields Medal, the Abel Prize has no age restriction. For Science Quiz,
    the Abel Prize is often asked through the pattern
    \[
        \text{year} \quad \longleftrightarrow \quad \text{laureate} \quad
        \longleftrightarrow \quad \text{famous theorem or area}.
    \]

    \vspace{1em}

    {\small
    \renewcommand{\arraystretch}{1.18}
    \begin{longtable}{@{}L{0.08\textwidth}L{0.32\textwidth}L{0.52\textwidth}@{}}
        \toprule
        Year & Laureate(s) & Quiz keywords \\
        \midrule
        \endfirsthead

        \toprule
        Year & Laureate(s) & Quiz keywords \\
        \midrule
        \endhead

        \midrule
        \multicolumn{3}{r}{Continued on the next page.} \\
        \endfoot

        \bottomrule
        \endlastfoot

        2003 & Jean-Pierre Serre
        & Algebraic topology, algebraic geometry, number theory; first Abel Prize laureate. \\

        2004 & Michael Atiyah; Isadore Singer
        & Atiyah--Singer index theorem. \\

        2005 & Peter D. Lax
        & Partial differential equations and applied mathematics. \\

        2006 & Lennart Carleson
        & Harmonic analysis; convergence of Fourier series. \\

        2007 & Srinivasa S. R. Varadhan
        & Probability theory; large deviations. \\

        2008 & John G. Thompson; Jacques Tits
        & Finite groups, classification of finite simple groups; buildings and group theory. \\

        2009 & Mikhail Gromov
        & Geometry, geometric group theory, symplectic geometry. \\

        2010 & John Tate
        & Number theory and arithmetic geometry. \\

        2011 & John Milnor
        & Topology, geometry, dynamics; exotic spheres. \\

        2012 & Endre Szemerédi
        & Discrete mathematics and theoretical computer science; Szemerédi's theorem. \\

        2013 & Pierre Deligne
        & Algebraic geometry; Weil conjectures. \\

        2014 & Yakov Sinai
        & Dynamical systems, ergodic theory, mathematical physics. \\

        2015 & John F. Nash Jr.; Louis Nirenberg
        & Nonlinear PDE, geometry, Nash embedding theorem. \\

        2016 & Andrew Wiles
        & Fermat's Last Theorem; modularity of elliptic curves. \\

        2017 & Yves Meyer
        & Wavelet theory; applications from signal compression to gravitational-wave analysis. \\

        2018 & Robert P. Langlands
        & Langlands program; representation theory and number theory. \\

        2019 & Karen Uhlenbeck
        & Geometric analysis and gauge theory; first woman Abel Prize laureate. \\

        2020 & Hillel Furstenberg; Gregory Margulis
        & Probability, dynamics, ergodic theory, Lie groups. \\

        2021 & László Lovász; Avi Wigderson
        & Discrete mathematics, theoretical computer science, complexity theory. \\

        2022 & Dennis P. Sullivan
        & Topology, geometric topology, dynamical systems. \\

        2023 & Luis A. Caffarelli
        & Nonlinear partial differential equations, free-boundary problems, Monge--Ampère equation. \\

        2024 & Michel Talagrand
        & Probability theory, functional analysis, spin glasses, concentration inequalities. \\

        2025 & Masaki Kashiwara
        & Algebraic analysis, $D$-modules, representation theory, crystal bases. \\

        2026 & Gerd Faltings
        & Arithmetic geometry; Mordell conjecture, Lang's conjectures, Diophantine geometry. \\
    \end{longtable}
    }

    \begin{fact}[Abel Prize landmarks]
        The following associations summarize the prize's main quiz anchors.
        \begin{itemize}
            \item Abel Prize $\leftrightarrow$ Norwegian Academy of Science and Letters.
            \item Abel Prize $\leftrightarrow$ Niels Henrik Abel.
            \item First Abel Prize laureate $\leftrightarrow$ Jean-Pierre Serre.
            \item First woman Abel Prize laureate $\leftrightarrow$ Karen Uhlenbeck.
            \item The Abel Prize money is on the order of one million US dollars.
            \item John Milnor $\leftrightarrow$ Fields Medal 1962 and Abel Prize 2011.
            \item Andrew Wiles $\leftrightarrow$ Abel Prize 2016 $\leftrightarrow$ Fermat's Last
            Theorem.
            \item Yves Meyer $\leftrightarrow$ Abel Prize 2017 $\leftrightarrow$ wavelets.
            \item Robert Langlands $\leftrightarrow$ Abel Prize 2018 $\leftrightarrow$ Langlands
            program.
            \item Gerd Faltings $\leftrightarrow$ Fields Medal 1986 and Abel Prize 2026.
        \end{itemize}
    \end{fact}

    \section*{Wolf Prize in Mathematics}

    The Wolf Prize is awarded by the Wolf Foundation. In the sciences, the Wolf
    Prize is awarded in Medicine, Agriculture, Mathematics, Chemistry, and
    Physics. It is especially important historically because many Wolf Prize in
    Mathematics laureates later received, or had already received, other major
    mathematical prizes. Unlike the Fields Medal, it is not restricted by age.

    \vspace{1em}

    {\small
    \renewcommand{\arraystretch}{1.16}
    \begin{longtable}{@{}L{0.09\textwidth}L{0.38\textwidth}L{0.45\textwidth}@{}}
        \toprule
        Year & Laureate(s) & Quiz keywords \\
        \midrule
        \endfirsthead

        \toprule
        Year & Laureate(s) & Quiz keywords \\
        \midrule
        \endhead

        \midrule
        \multicolumn{3}{r}{Continued on the next page.} \\
        \endfoot

        \bottomrule
        \endlastfoot

        1978 & Israel Gelfand; Carl L. Siegel
        & Representation theory, functional analysis; number theory, several complex variables. \\

        1979 & Jean Leray; André Weil
        & PDE and algebraic topology; algebraic geometry and number theory. \\

        1980 & Henri Cartan; Andrey Kolmogorov
        & Several complex variables and algebraic topology; probability theory and dynamical systems. \\

        1981 & Lars Ahlfors; Oscar Zariski
        & Complex analysis and Riemann surfaces; algebraic geometry. \\

        1982 & Mark Krein; Hassler Whitney
        & Functional analysis; differential topology and Whitney embedding. \\

        1983 & Shiing-Shen Chern
        & Global differential geometry, Chern classes. \\

        1984 & Paul Erd\H{o}s
        & Combinatorics, number theory, probabilistic method. \\

        1985 & Kunihiko Kodaira; Hans Lewy
        & Complex geometry; partial differential equations. \\

        1986 & Samuel Eilenberg; Atle Selberg
        & Algebraic topology and category theory; analytic number theory and Selberg trace formula. \\

        1987 & Kiyoshi Itô; Peter D. Lax
        & Stochastic calculus; partial differential equations and applied mathematics. \\

        1988 & Lars Hörmander; Friedrich Hirzebruch
        & Partial differential equations; topology, characteristic classes, Hirzebruch--Riemann--Roch. \\

        1989 & Alberto Calderón; John Milnor
        & Harmonic analysis and singular integrals; topology and exotic spheres. \\

        1990 & Ennio De Giorgi; Ilya Piatetski-Shapiro
        & Minimal surfaces and regularity theory; automorphic forms. \\

        1992 & Lennart Carleson; John G. Thompson
        & Harmonic analysis; finite group theory. \\

        1993 & Mikhail Gromov; Jacques Tits
        & Geometry and geometric group theory; buildings and group theory. \\

        1995/96 & Robert Langlands; Andrew Wiles
        & Langlands program; Fermat's Last Theorem. \\

        1996/97 & Joseph Keller; Yakov Sinai
        & Applied mathematics and waves; dynamical systems and mathematical physics. \\

        1999 & László Lovász; Elias M. Stein
        & Discrete mathematics and algorithms; harmonic analysis. \\

        2000 & Raoul Bott; Jean-Pierre Serre
        & Topology and geometry; algebraic topology, algebraic geometry, number theory. \\

        2001 & Vladimir Arnold; Saharon Shelah
        & Dynamical systems and singularity theory; mathematical logic and set theory. \\

        2002/03 & Mikio Sato; John Tate
        & Algebraic analysis and hyperfunctions; arithmetic geometry. \\

        2005 & Gregory Margulis; Serge Novikov
        & Lie groups, ergodic theory, discrete subgroups; topology and geometry. \\

        2006/07 & Stephen Smale; Hillel Furstenberg
        & Differential topology and dynamical systems; ergodic theory and probability. \\

        2008/09 & Pierre Deligne; Phillip Griffiths; David Mumford
        & Algebraic geometry, Hodge theory, moduli spaces. \\

        2010 & Shing-Tung Yau; Dennis Sullivan
        & Geometric analysis; topology and dynamics. \\

        2012 & Michael Aschbacher; Luis Caffarelli
        & Finite groups; nonlinear PDE and free-boundary problems. \\

        2013 & George Mostow; Michael Artin
        & Rigidity, geometry, Lie groups; algebraic geometry. \\

        2014 & Peter Sarnak
        & Number theory, analysis, geometry, automorphic forms. \\

        2015 & James Arthur
        & Trace formula, automorphic representations. \\

        2017 & Richard Schoen; Charles Fefferman
        & Geometric analysis; analysis and partial differential equations. \\

        2018 & Alexander Beilinson; Vladimir Drinfeld
        & Geometric representation theory, mathematical physics, Langlands program. \\

        2019 & Jean-François Le Gall; Gregory Lawler
        & Probability theory, random walks, Brownian motion, SLE. \\

        2020 & Simon Donaldson; Yakov Eliashberg
        & Differential geometry, topology, gauge theory; symplectic and contact topology. \\

        2022 & George Lusztig
        & Representation theory and related areas. \\

        2023 & Ingrid Daubechies
        & Wavelet theory, time-frequency analysis, signal processing. \\

        2024 & Adi Shamir; Noga Alon
        & Mathematical cryptography; combinatorics, graph theory, theoretical computer science. \\
    \end{longtable}
    }

    \begin{fact}[Wolf Prize landmarks]
        The following associations summarize the prize's main quiz anchors.
        \begin{itemize}
            \item Wolf Prize $\leftrightarrow$ Wolf Foundation.
            \item Wolf Prize in Mathematics $\leftrightarrow$ no Fields-style age restriction.
            \item Paul Erd\H{o}s $\leftrightarrow$ Wolf Prize 1984 $\leftrightarrow$ combinatorics and
            probabilistic method.
            \item John Milnor $\leftrightarrow$ Wolf Prize 1989; together with the Fields Medal and
            Abel Prize, this makes him a standard example connected to all three major mathematics
            prizes.
            \item Andrew Wiles and Robert Langlands shared the 1995/96 Wolf Prize in Mathematics.
            \item Ingrid Daubechies $\leftrightarrow$ Wolf Prize 2023 $\leftrightarrow$ wavelets.
            \item Adi Shamir and Noga Alon $\leftrightarrow$ Wolf Prize 2024 $\leftrightarrow$
            cryptography and combinatorics.
        \end{itemize}
    \end{fact}

    \begin{fact}[Prize recognition]
        For quick Science Quiz recall, remember the following map first.
        \begin{itemize}
            \item Fields Medal $\leftrightarrow$ ICM $\leftrightarrow$ IMU $\leftrightarrow$ every
            four years $\leftrightarrow$ age condition.
            \item Abel Prize $\leftrightarrow$ Norwegian Academy of Science and Letters
            $\leftrightarrow$ no age restriction.
            \item Wolf Prize $\leftrightarrow$ Wolf Foundation $\leftrightarrow$ no Fields-style age
            restriction.
            \item Fields Medal questions often ask for the year, ICM location, laureate, and
            keyword.
            \item Abel and Wolf questions often ask for the laureate--area association.
        \end{itemize}
    \end{fact}

    \chapter{Mathematicians, History, and Anecdotes}

    Science Quiz repeatedly asks for a mathematician from a birth year, a named
    theorem, a field, a book, or a short biographical clue. The most efficient
    preparation is therefore chronological: recognize the person, place the
    person in the correct period, and attach one or two unmistakable
    achievements.

    \section*{Mathematicians}

    The following table is ordered by birth. Ancient dates are approximate, and
    ``fl.'' indicates the period in which a person is known to have worked.
    Living mathematicians are marked by their year of birth. A final
    contemporary entry whose birth year is not publicly listed is separated
    from the chronological sequence.

    \vspace{1em}

    {\scriptsize
    \renewcommand{\arraystretch}{1.13}
    \begin{longtable}{@{}L{0.20\textwidth}L{0.13\textwidth}L{0.18\textwidth}L{0.41\textwidth}@{}}
        \toprule
        Mathematician & Life & Main fields & Quiz anchors \\
        \midrule
        \endfirsthead

        \toprule
        Mathematician & Life & Main fields & Quiz anchors \\
        \midrule
        \endhead

        \midrule
        \multicolumn{4}{r}{Continued on the next page.} \\
        \endfoot

        \bottomrule
        \endlastfoot

        Thales of Miletus
        & c. 624--c. 546 BCE
        & Geometry
        & Early deductive geometry; Thales' theorem; a right angle in a
          semicircle. \\

        Pythagoras
        & c. 570--c. 495 BCE
        & Geometry, number theory
        & Pythagorean theorem; Pythagorean school; numerical ratios in music. \\

        Euclid
        & fl. c. 300 BCE
        & Geometry, foundations
        & \textit{Elements}; axiomatic method; Euclidean algorithm; infinitely
          many primes. \\

        Archimedes
        & c. 287--212 BCE
        & Geometry, mechanics
        & Method of exhaustion; bounds for $\pi$; lever; buoyancy; sphere and
          circumscribed cylinder. \\

        Apollonius of Perga
        & c. 240--c. 190 BCE
        & Geometry
        & \textit{Conics}; systematic study of ellipses, parabolas, and
          hyperbolas. \\

        Diophantus
        & c. 200--c. 284
        & Number theory, algebra
        & \textit{Arithmetica}; Diophantine equations. \\

        Liu Hui
        & c. 225--c. 295
        & Arithmetic, geometry
        & Commentary on \textit{The Nine Chapters}; polygon method for $\pi$;
          elimination procedures. \\

        Hypatia
        & c. 355--415
        & Geometry, astronomy
        & Alexandrian teacher and commentator; associated with editions of
          Diophantus and Apollonius. \\

        Zu Chongzhi
        & 429--500
        & Numerical mathematics, astronomy
        & $3.1415926<\pi<3.1415927$; the approximation
          $\pi\approx 355/113$. \\

        Aryabhata
        & 476--550
        & Arithmetic, trigonometry, astronomy
        & \textit{Aryabhatiya}; sine tables; place-value computation and
          astronomical algorithms. \\

        Brahmagupta
        & 598--c. 668
        & Algebra, number theory
        & Rules for zero and negative numbers; Brahmagupta's formula; quadratic
          equations. \\

        Muhammad al-Khwarizmi
        & c. 780--c. 850
        & Algebra, arithmetic
        & Systematic algebra; the words \emph{algebra} and \emph{algorithm}. \\

        Omar Khayyam
        & 1048--1131
        & Algebra, geometry
        & Classification of cubic equations and geometric solutions by conic
          sections. \\

        Bhāskara II
        & 1114--1185
        & Algebra, number theory, astronomy
        & \textit{Līlāvatī}; indeterminate equations; cyclic methods related to
          Pell equations. \\

        Fibonacci
        & c. 1170--c. 1240
        & Arithmetic, number theory
        & \textit{Liber Abaci}; spread of Hindu--Arabic numerals in Europe;
          Fibonacci sequence. \\

        Madhava of Sangamagrama
        & c. 1340--c. 1425
        & Infinite series, trigonometry
        & Kerala school; series for $\pi$, sine, and cosine. \\

        Jamshid al-Kashi
        & c. 1380--1429
        & Numerical mathematics, astronomy
        & Decimal fractions; high-precision computation of $\pi$; law of
          cosines. \\

        Niccolò Tartaglia
        & c. 1499--1557
        & Algebra
        & A method for depressed cubic equations; the Cardano--Tartaglia
          priority dispute. \\

        Gerolamo Cardano
        & 1501--1576
        & Algebra, probability
        & \textit{Ars Magna}; cubic and quartic equations; early probability. \\

        Lodovico Ferrari
        & 1522--1565
        & Algebra
        & General solution of the quartic equation, published in Cardano's
          \textit{Ars Magna}. \\

        Rafael Bombelli
        & 1526--1572
        & Algebra
        & \textit{L'Algebra}; systematic rules for complex numbers in solutions
          of cubic equations. \\

        François Viète
        & 1540--1603
        & Algebra, trigonometry
        & Systematic symbolic algebra; Viète's formulas. \\

        Simon Stevin
        & 1548--1620
        & Arithmetic, mechanics
        & Promoted decimal fractions in Europe; work on statics and hydrostatics. \\

        John Napier
        & 1550--1617
        & Computation, logarithms
        & Introduced logarithms in print in 1614; Napier's bones. \\

        Galileo Galilei
        & 1564--1642
        & Mathematical physics
        & Mathematical description of motion; telescopic astronomy. \\

        Johannes Kepler
        & 1571--1630
        & Geometry, astronomy
        & Laws of planetary motion; Kepler conjecture on sphere packing. \\

        Marin Mersenne
        & 1588--1648
        & Number theory, acoustics
        & Mersenne primes $2^p-1$; scientific correspondence network in
          seventeenth-century Europe. \\

        Girard Desargues
        & 1591--1661
        & Projective geometry
        & Desargues' theorem; foundational ideas of projective geometry. \\

        René Descartes
        & 1596--1650
        & Analytic geometry, algebra
        & Cartesian coordinates; \textit{La Géométrie}; rule of signs. \\

        Bonaventura Cavalieri
        & 1598--1647
        & Geometry
        & Method of indivisibles, an important precursor to integral calculus. \\

        Pierre de Fermat
        & 1607--1665
        & Number theory, geometry, probability
        & Fermat's Last Theorem; little theorem; analytic geometry; probability
          correspondence with Pascal. \\

        John Wallis
        & 1616--1703
        & Analysis, geometry
        & Introduced the symbol $\infty$; Wallis product for $\pi$; infinitesimal
          methods. \\

        Blaise Pascal
        & 1623--1662
        & Probability, geometry
        & Pascal's triangle; foundations of probability with Fermat; projective
          geometry. \\

        Christiaan Huygens
        & 1629--1695
        & Probability, mechanics
        & First printed treatise on probability; pendulum and wave theory. \\

        Isaac Barrow
        & 1630--1677
        & Geometry, analysis
        & Geometric form of the fundamental theorem of calculus; Newton's
          teacher at Cambridge. \\

        Isaac Newton
        & 1643--1727
        & Calculus, mechanics
        & Fluxions; fundamental theorem of calculus; generalized binomial
          theorem; laws of motion and gravitation. \\

        Gottfried Wilhelm Leibniz
        & 1646--1716
        & Calculus, logic
        & Independent development of calculus; $\dd x$ and $\int$ notation;
          binary arithmetic. \\

        Jakob Bernoulli
        & 1655--1705
        & Probability, analysis
        & Bernoulli numbers; law of large numbers; \textit{Ars Conjectandi}. \\

        Guillaume de l'Hôpital
        & 1661--1704
        & Analysis
        & First calculus textbook; l'Hôpital's rule, developed from lessons by
          Johann Bernoulli. \\

        Abraham de Moivre
        & 1667--1754
        & Probability, complex numbers
        & De Moivre's formula; normal approximation to the binomial
          distribution. \\

        Johann Bernoulli
        & 1667--1748
        & Analysis, mechanics
        & Brachistochrone problem; early calculus; teacher of Euler and
          l'Hôpital. \\

        Brook Taylor
        & 1685--1731
        & Analysis
        & Taylor's theorem and Taylor series; finite differences. \\

        Christian Goldbach
        & 1690--1764
        & Number theory
        & Goldbach conjecture, proposed in correspondence with Euler. \\

        James Stirling
        & 1692--1770
        & Analysis, combinatorics
        & Stirling's approximation; Stirling numbers. \\

        Colin Maclaurin
        & 1698--1746
        & Analysis, geometry
        & Maclaurin series; \textit{Treatise of Fluxions}. \\

        Daniel Bernoulli
        & 1700--1782
        & Probability, fluid mechanics
        & Bernoulli principle; expected utility and the St.\ Petersburg
          paradox. \\

        Thomas Bayes
        & c. 1701--1761
        & Probability
        & Bayes' theorem; posthumous 1763 publication. \\

        Leonhard Euler
        & 1707--1783
        & Analysis, number theory, geometry
        & Basel problem; Euler formula and identity; graph theory; polyhedron
          formula; standardized much modern notation. \\

        Jean le Rond d'Alembert
        & 1717--1783
        & Analysis, mechanics
        & Wave equation; d'Alembert's principle; work on the fundamental theorem
          of algebra. \\

        Maria Gaetana Agnesi
        & 1718--1799
        & Analysis
        & \textit{Instituzioni analitiche}; one of the earliest comprehensive
          calculus textbooks; the witch of Agnesi curve. \\

        Joseph-Louis Lagrange
        & 1736--1813
        & Analysis, mechanics, number theory
        & Lagrangian mechanics; multipliers; four-square theorem. \\

        Gaspard Monge
        & 1746--1818
        & Geometry
        & Descriptive geometry; Monge formulation of optimal transport. \\

        Pierre-Simon Laplace
        & 1749--1827
        & Probability, celestial mechanics
        & Laplace transform; Laplacian; analytic theory of probability. \\

        Adrien-Marie Legendre
        & 1752--1833
        & Number theory, analysis
        & Legendre symbol; elliptic integrals; least squares. \\

        Joseph Fourier
        & 1768--1830
        & Analysis, PDE
        & Fourier series and transform; heat equation. \\

        Sophie Germain
        & 1776--1831
        & Number theory, elasticity
        & Work on Fermat's Last Theorem and vibrating plates; used the
          pseudonym M.\ LeBlanc. \\

        Carl Friedrich Gauss
        & 1777--1855
        & Number theory, geometry, statistics
        & \textit{Disquisitiones Arithmeticae}; quadratic reciprocity; Gaussian
          distribution; least squares; curvature. \\

        Siméon Denis Poisson
        & 1781--1840
        & Analysis, probability, physics
        & Poisson distribution and equation; mathematical physics and potential
          theory. \\

        Bernard Bolzano
        & 1781--1848
        & Analysis, logic
        & Early rigorous analysis; intermediate value theorem; Bolzano--
          Weierstrass theorem. \\

        Augustin-Louis Cauchy
        & 1789--1857
        & Analysis
        & Rigorous limits and convergence; Cauchy sequences; integral theorem;
          determinant theory. \\

        August Ferdinand Möbius
        & 1790--1868
        & Geometry, number theory
        & Möbius strip; Möbius transformations and inversion function. \\

        Charles Babbage
        & 1791--1871
        & Computation, numerical mathematics
        & Difference Engine and Analytical Engine. \\

        Nikolai Lobachevsky
        & 1792--1856
        & Geometry
        & Independent creation of hyperbolic, non-Euclidean geometry. \\

        George Green
        & 1793--1841
        & Analysis, mathematical physics
        & Green's theorem, Green's identities, and Green functions. \\

        Niels Henrik Abel
        & 1802--1829
        & Algebra, analysis
        & Impossibility of a radical formula for the general quintic; elliptic
          functions; Abel Prize namesake. \\

        János Bolyai
        & 1802--1860
        & Geometry
        & Independent creation of hyperbolic geometry. \\

        Carl Gustav Jacobi
        & 1804--1851
        & Analysis, mechanics, number theory
        & Jacobian determinant; elliptic functions; Hamilton--Jacobi theory. \\

        Peter Gustav Lejeune Dirichlet
        & 1805--1859
        & Number theory, analysis
        & Primes in arithmetic progressions; Dirichlet characters; modern
          function concept. \\

        William Rowan Hamilton
        & 1805--1865
        & Algebra, mechanics
        & Quaternions; Hamiltonian mechanics; Hamilton--Jacobi equation. \\

        Augustus De Morgan
        & 1806--1871
        & Logic, algebra
        & De Morgan's laws; formal logic; mathematical induction terminology. \\

        Joseph Liouville
        & 1809--1882
        & Analysis, number theory
        & First proof that transcendental numbers exist; Liouville's theorem. \\

        Hermann Grassmann
        & 1809--1877
        & Algebra, geometry
        & \textit{Ausdehnungslehre}; vector spaces and exterior algebra. \\

        Ernst Kummer
        & 1810--1893
        & Number theory, algebra
        & Ideal numbers; regular primes; major progress on Fermat's Last
          Theorem. \\

        Évariste Galois
        & 1811--1832
        & Algebra
        & Galois theory; solvability by radicals through groups. \\

        James Joseph Sylvester
        & 1814--1897
        & Algebra, number theory
        & Matrix and invariant theory; coined the term ``matrix''; Sylvester's
          law of inertia. \\

        Karl Weierstrass
        & 1815--1897
        & Analysis
        & $\eps$--$\delta$ rigor; uniform convergence; nowhere-differentiable
          function. \\

        George Boole
        & 1815--1864
        & Logic, algebra
        & Boolean algebra; algebraic treatment of logic. \\

        Ada Lovelace
        & 1815--1852
        & Computation
        & Notes on Babbage's Analytical Engine; an algorithm for Bernoulli
          numbers often described as the first published computer program. \\

        George Gabriel Stokes
        & 1819--1903
        & Analysis, mathematical physics
        & Stokes' theorem; fluid equations; optics. \\

        Pafnuty Chebyshev
        & 1821--1894
        & Number theory, approximation, probability
        & Chebyshev polynomials and inequality; work on prime distribution. \\

        Arthur Cayley
        & 1821--1895
        & Algebra, geometry
        & Matrix algebra; Cayley's theorem; Cayley--Hamilton theorem. \\

        Charles Hermite
        & 1822--1901
        & Analysis, number theory
        & Proved $e$ transcendental in 1873; Hermite polynomials and normal
          forms. \\

        Leopold Kronecker
        & 1823--1891
        & Algebra, number theory
        & Kronecker product and delta; arithmetic approach to algebra. \\

        Bernhard Riemann
        & 1826--1866
        & Analysis, geometry, number theory
        & Riemann surfaces and integral; Riemannian geometry; zeta function and
          hypothesis. \\

        Richard Dedekind
        & 1831--1916
        & Algebra, foundations, number theory
        & Dedekind cuts; ideals; rigorous construction of the real numbers. \\

        Camille Jordan
        & 1838--1922
        & Algebra, analysis
        & Jordan normal form; Jordan curve theorem; permutation groups. \\

        Sophus Lie
        & 1842--1899
        & Geometry, differential equations
        & Lie groups and Lie algebras; continuous symmetries. \\

        Georg Cantor
        & 1845--1918
        & Set theory
        & Founder of set theory; diagonal argument; cardinalities; Cantor set. \\

        Wilhelm Killing
        & 1847--1923
        & Lie theory, geometry
        & Classification program for simple Lie algebras; Killing form. \\

        Gottlob Frege
        & 1848--1925
        & Logic, foundations
        & Predicate logic; \textit{Begriffsschrift}; logicist foundations of
          arithmetic. \\

        Felix Klein
        & 1849--1925
        & Geometry, group theory
        & Erlangen program; Klein bottle; Klein four-group. \\

        Ferdinand Georg Frobenius
        & 1849--1917
        & Algebra, analysis
        & Frobenius theorem and endomorphism; character theory of finite groups. \\

        Sofia Kovalevskaya
        & 1850--1891
        & Analysis, PDE
        & Cauchy--Kovalevskaya theorem; rigid-body dynamics. \\

        Ferdinand von Lindemann
        & 1852--1939
        & Number theory, analysis
        & Proved $\pi$ transcendental in 1882, establishing the impossibility
          of squaring the circle. \\

        Henri Poincaré
        & 1854--1912
        & Topology, dynamics, geometry
        & Poincaré conjecture; qualitative dynamics; foundational topology. \\

        Andrey Markov
        & 1856--1922
        & Probability, number theory
        & Markov chains; Markov inequality; stochastic dependence. \\

        Giuseppe Peano
        & 1858--1932
        & Logic, foundations
        & Peano axioms; symbolic logic; space-filling Peano curve. \\

        David Hilbert
        & 1862--1943
        & Foundations, algebra, geometry
        & Twenty-three problems; Hilbert spaces and basis theorem; axiomatic
          geometry. \\

        John Charles Fields
        & 1863--1932
        & Complex analysis
        & Canadian mathematician who proposed and endowed the Fields Medal. \\

        Hermann Minkowski
        & 1864--1909
        & Geometry, number theory
        & Geometry of numbers; Minkowski spacetime; convex body theorem. \\

        Jacques Hadamard
        & 8 Dec 1865--17 Oct 1963
        & Analysis, number theory
        & Prime number theorem; Cauchy--Hadamard theorem; Hadamard matrices. \\

        Charles-Jean de la Vallée Poussin
        & 1866--1962
        & Analysis, number theory
        & Independent proof of the prime number theorem in 1896. \\

        Felix Hausdorff
        & 1868--1942
        & Topology, set theory
        & Hausdorff spaces; Hausdorff measure and dimension. \\

        Élie Cartan
        & 1869--1951
        & Differential geometry, Lie theory
        & Lie groups and algebras; differential forms; moving frames. \\

        Ernst Zermelo
        & 1871--1953
        & Set theory, foundations
        & Axiomatization of set theory; well-ordering theorem; Zermelo--Fraenkel
          set theory. \\

        Bertrand Russell
        & 1872--1970
        & Logic, foundations
        & Russell's paradox; \textit{Principia Mathematica} with Whitehead. \\

        Alfred Young
        & 1873--1940
        & Algebra, representation theory
        & Young diagrams and Young tableaux for symmetric groups. \\

        Henri Lebesgue
        & 1875--1941
        & Analysis, measure theory
        & Lebesgue measure and integral; dominated convergence theorem. \\

        G. H. Hardy
        & 1877--1947
        & Analysis, number theory
        & Hardy--Weinberg principle; Hardy--Littlewood circle method;
          collaboration with Ramanujan. \\

        Maurice Fréchet
        & 1878--1973
        & Analysis, topology
        & Metric spaces; Fréchet derivative; abstract functional spaces. \\

        Frigyes Riesz
        & 1880--1956
        & Functional analysis
        & Riesz representation theorem; $L^p$ spaces and operator theory. \\

        L. E. J. Brouwer
        & 1881--1966
        & Topology, foundations
        & Brouwer fixed-point theorem; invariance of domain; intuitionism. \\

        Wacław Sierpiński
        & 1882--1969
        & Set theory, topology
        & Sierpiński triangle and carpet; continuum and descriptive set theory. \\

        Emmy Noether
        & 1882--1935
        & Algebra, mathematical physics
        & Noetherian rings and modules; symmetry--conservation theorem. \\

        J. E. Littlewood
        & 1885--1977
        & Analysis, number theory
        & Hardy--Littlewood circle method and conjectures; analytic number
          theory. \\

        Hermann Weyl
        & 1885--1955
        & Geometry, representation theory, physics
        & Weyl groups and character formula; gauge ideas; foundations. \\

        Hugo Steinhaus
        & 1887--1972
        & Analysis, probability
        & Lwów school; Steinhaus theorem; co-founder of \textit{Studia
          Mathematica}. \\

        George Pólya
        & 1887--1985
        & Combinatorics, probability
        & Pólya enumeration theorem; random walks; \textit{How to Solve It}. \\

        Srinivasa Ramanujan
        & 1887--1920
        & Number theory, analysis
        & Partitions, modular forms, mock theta functions; $1729$. \\

        Richard Courant
        & 1888--1972
        & Analysis, PDE
        & Courant Institute namesake; finite elements; Courant--Friedrichs--Lewy
          condition. \\

        Ludwig Wittgenstein
        & 1889--1951
        & Logic, philosophy
        & \textit{Tractatus}; language games; major influence on logical
          positivism and philosophy of language. \\

        Stefan Banach
        & 1892--1945
        & Functional analysis
        & Banach spaces; contraction mapping theorem; Lwów school and
          \textit{The Scottish Book}. \\

        Norbert Wiener
        & 1894--1964
        & Analysis, probability
        & Wiener process and measure; harmonic analysis; founder of
          cybernetics. \\

        Emil Artin
        & 1898--1962
        & Algebra, number theory
        & Artin reciprocity; braid groups; modern algebraic number theory. \\

        Oscar Zariski
        & 1899--1986
        & Algebraic geometry
        & Zariski topology; foundations of modern algebraic geometry. \\

        Juliusz Schauder
        & 1899--1943
        & Functional analysis, PDE
        & Schauder fixed-point theorem and Schauder estimates. \\

        Antoni Zygmund
        & 1900--1992
        & Harmonic analysis
        & \textit{Trigonometric Series}; Chicago school of harmonic analysis. \\

        Alfred Tarski
        & 1901--1983
        & Logic, foundations
        & Semantic definition of truth; Tarski's theorem; Banach--Tarski
          paradox. \\

        Andrey Kolmogorov
        & 1903--1987
        & Probability, dynamics
        & Axioms of probability; Kolmogorov complexity; turbulence and KAM
          theory. \\

        Alonzo Church
        & 1903--1995
        & Logic, computation
        & Lambda calculus; Church's theorem; Church--Turing thesis. \\

        John von Neumann
        & 1903--1957
        & Analysis, computation, game theory
        & Operator algebras; quantum mechanics; minimax theorem; stored-program
          architecture. \\

        Stanisław Mazur
        & 1905--1981
        & Functional analysis
        & Lwów school; Banach--Mazur theorem; Scottish Book basis problem and
          its live-goose prize. \\

        Kurt Gödel
        & 1906--1978
        & Logic, foundations
        & Completeness and incompleteness theorems; consistency of choice and
          CH with ZF. \\

        André Weil
        & 1906--1998
        & Number theory, algebraic geometry
        & Weil conjectures; adeles; founding member of Bourbaki. \\

        Hans Fitting
        & 1906--1938
        & Group theory, algebra
        & Fitting subgroup, lemma, decomposition, and ideals. \\

        Lars Ahlfors
        & 1907--1996
        & Complex analysis
        & Riemann surfaces; Ahlfors theory; one of the first two Fields
          Medalists. \\

        Claude Chevalley
        & 1909--1984
        & Algebra, number theory
        & Chevalley groups; class field theory; founding member of Bourbaki. \\

        Stanisław Ulam
        & 1909--1984
        & Set theory, probability
        & Monte Carlo method; Ulam spiral; Scottish Book and Manhattan Project. \\

        Saunders Mac Lane
        & 1909--2005
        & Algebra, topology
        & Co-founder of category theory; Eilenberg--Mac Lane spaces. \\

        Shiing-Shen Chern
        & 1911--2004
        & Differential geometry, topology
        & Chern classes; Chern--Gauss--Bonnet theorem; Chern Medal namesake. \\

        Alan Turing
        & 1912--1954
        & Logic, computation
        & Turing machine; computability; codebreaking; Turing test. \\

        Paul Erd\H{o}s
        & 1913--1996
        & Number theory, combinatorics
        & Probabilistic method; Erd\H{o}s number; prolific collaboration. \\

        Israel Gelfand
        & 1913--2009
        & Functional analysis, representation theory
        & Gelfand representation; representation theory; influential seminar. \\

        Samuel Eilenberg
        & 1913--1998
        & Algebraic topology
        & Co-founder of category theory and homological algebra;
          Eilenberg--Mac Lane spaces. \\

        Mark Kac
        & 1914--1984
        & Probability, mathematical physics
        & Feynman--Kac formula; ``Can one hear the shape of a drum?'' \\

        George Dantzig
        & 1914--2005
        & Optimization, statistics
        & Linear programming; simplex method; the mistaken-homework anecdote. \\

        Laurent Schwartz
        & 1915--2002
        & Analysis
        & Theory of distributions; Fields Medal 1950. \\

        Claude Shannon
        & 1916--2001
        & Information theory
        & Entropy, channel capacity, and the mathematical theory of
          communication. \\

        Atle Selberg
        & 1917--2007
        & Analytic number theory
        & Selberg sieve and trace formula; elementary proof of the prime number
          theorem with Erd\H{o}s. \\

        Benoit Mandelbrot
        & 1924--2010
        & Geometry, probability
        & Fractal geometry; Mandelbrot set; popularized the term ``fractal.'' \\

        Louis Nirenberg
        & 1925--2020
        & Analysis, PDE
        & Fundamental estimates for elliptic PDE; Abel Prize 2015 with John
          Nash. \\

        Peter D. Lax
        & born 1926
        & Analysis, PDE
        & Lax pair and Lax equivalence theorem; numerical analysis; Abel Prize
          2005. \\

        Jean-Pierre Serre
        & born 1926
        & Algebraic topology, geometry, number theory
        & Youngest Fields Medalist; first Abel laureate; Serre duality and
          spectral sequences. \\

        Alexander Grothendieck
        & 1928--2014
        & Algebraic geometry
        & Schemes, topoi, étale cohomology; transformed modern algebraic
          geometry. \\

        John Nash
        & 1928--2015
        & Geometry, game theory, PDE
        & Nash equilibrium; Nash embedding theorem; Abel Prize 2015. \\

        Michael Atiyah
        & 1929--2019
        & Geometry, topology
        & Atiyah--Singer index theorem; K-theory; Fields and Abel prizes. \\

        Stephen Smale
        & born 1930
        & Topology, dynamics
        & Higher-dimensional Poincaré conjecture; Smale horseshoe; Smale's
          problems. \\

        Lars Hörmander
        & 1931--2012
        & Analysis, PDE
        & Linear partial differential operators; microlocal analysis; Fields
          Medal 1962. \\

        John Milnor
        & born 1931
        & Topology, geometry, dynamics
        & Exotic $7$-spheres; Milnor fibration; Fields, Wolf, and Abel prizes. \\

        Paul Cohen
        & 1934--2007
        & Logic, set theory
        & Forcing; independence of the continuum hypothesis and axiom of
          choice. \\

        Yakov Sinai
        & born 1935
        & Dynamical systems, probability
        & Sinai billiards; Kolmogorov--Sinai entropy; Abel Prize 2014. \\

        Robert Langlands
        & born 1936
        & Number theory, representation theory
        & Langlands program; reciprocity and functoriality; Abel Prize 2018. \\

        Vladimir Arnold
        & 1937--2010
        & Dynamical systems, geometry
        & KAM theory; Arnold conjecture; singularity theory. \\

        John Horton Conway
        & 1937--2020
        & Algebra, combinatorics
        & Surreal numbers; cellular automaton \textit{Life}; Conway groups and
          knot notation. \\

        Donald Knuth
        & born 1938
        & Algorithms, combinatorics
        & \textit{The Art of Computer Programming}; TeX; analysis of
          algorithms. \\

        Alan Baker
        & 1939--2018
        & Number theory
        & Linear forms in logarithms; transcendence theory; Fields Medal 1970. \\

        Karen Uhlenbeck
        & born 1942
        & Geometric analysis
        & Gauge theory and harmonic maps; first woman Abel laureate. \\

        Mikhail Gromov
        & born 1943
        & Geometry, topology
        & Gromov--Hausdorff distance; geometric group theory; Abel Prize 2009. \\

        Per Enflo
        & born 1944
        & Functional analysis
        & Solved Mazur's Scottish Book basis problem and received the promised
          live goose. \\

        Pierre Deligne
        & born 1944
        & Algebraic geometry
        & Proof of the Weil conjectures; Fields, Wolf, and Abel prizes. \\

        Grigory Margulis
        & born 1946
        & Lie groups, dynamics
        & Superrigidity, arithmeticity, and homogeneous dynamics; Fields and
          Abel prizes. \\

        William Thurston
        & 1946--2012
        & Topology, geometry
        & Geometrization conjecture; eight geometries; low-dimensional
          topology. \\

        Shing-Tung Yau
        & born 1949
        & Geometric analysis
        & Calabi conjecture; Ricci-flat metrics; Fields Medal 1982. \\

        Edward Witten
        & born 1951
        & Mathematical physics
        & First physicist to receive the Fields Medal; quantum field theory,
          topology, and string theory. \\

        Andrew Wiles
        & born 1953
        & Number theory
        & Fermat's Last Theorem via modularity; special IMU plaque and Abel
          Prize. \\

        Gerd Faltings
        & born 1954
        & Arithmetic geometry
        & Mordell conjecture; Fields Medal 1986; Abel Prize 2026. \\

        Ingrid Daubechies
        & born 1954
        & Analysis, applied mathematics
        & Orthogonal wavelets; Daubechies wavelets; Wolf Prize 2023. \\

        Efim Zelmanov
        & born 1955
        & Algebra, group theory
        & Restricted Burnside problem; Fields Medal 1994. \\

        Simon Donaldson
        & born 1957
        & Geometry, topology
        & Gauge theory and smooth $4$-manifolds; Fields Medal 1986. \\

        Maxim Kontsevich
        & born 1964
        & Geometry, mathematical physics
        & Deformation quantization; mirror symmetry; Fields Medal 1998. \\

        Grigori Perelman
        & born 1966
        & Geometry, topology
        & Poincaré and geometrization conjectures via Ricci flow; declined the
          Fields Medal and Millennium Prize. \\

        Elon Lindenstrauss
        & born 1970
        & Ergodic theory, number theory
        & Measure rigidity and applications to number theory; Fields Medal
          2010. \\

        Cédric Villani
        & born 1973
        & PDE, optimal transport
        & Boltzmann equation and Landau damping; Fields Medal 2010;
          \textit{Théorème vivant}. \\

        Manjul Bhargava
        & born 1974
        & Number theory
        & Higher composition laws; arithmetic statistics; Fields Medal 2014. \\

        Terence Tao
        & born 1975
        & Analysis, PDE, combinatorics
        & Green--Tao theorem; harmonic analysis; Fields Medal 2006. \\

        Martin Hairer
        & born 1975
        & Probability, stochastic PDE
        & Theory of regularity structures; Fields Medal 2014. \\

        Maryam Mirzakhani
        & 1977--2017
        & Geometry, dynamics
        & Moduli spaces and Riemann surfaces; first woman Fields Medalist. \\

        Caucher Birkar
        & born 1978
        & Algebraic geometry
        & Birational geometry and boundedness of Fano varieties; Fields Medal
          2018. \\

        Akshay Venkatesh
        & born 1981
        & Number theory, dynamics
        & Number theory and homogeneous dynamics; Fields Medal 2018. \\

        June Huh
        & born 1983
        & Combinatorics, algebraic geometry
        & Hodge theory for matroids; log-concavity conjectures; Fields Medal
          2022. \\

        Alessio Figalli
        & born 1984
        & Analysis, optimal transport
        & Optimal transport and PDE; Fields Medal 2018. \\

        Maryna Viazovska
        & born 1984
        & Number theory, discrete geometry
        & Optimal sphere packing in dimensions $8$ and $24$; Fields Medal
          2022. \\

        Hugo Duminil-Copin
        & born 1985
        & Probability, mathematical physics
        & Phase transitions and percolation; Fields Medal 2022. \\

        James Maynard
        & born 1987
        & Number theory
        & Small gaps between primes and Diophantine approximation; Fields Medal
          2022. \\

        Peter Scholze
        & born 1987
        & Arithmetic geometry
        & Perfectoid spaces and $p$-adic geometry; Fields Medal 2018. \\

        \midrule
        Botong Wang
        & active in the 21st century
        & Algebraic geometry, topology
        & Topology of algebraic varieties; interactions with combinatorics and
          applied mathematics; his birth year is not publicly listed. \\
    \end{longtable}
    }

    \section*{Mathematical History}

    The following chronology records changes that altered mathematical
    language, method, or subject matter. It is intentionally written as short
    recall statements rather than as a continuous historical essay.

    \begin{itemize}
        \item \textbf{c. 1800--1600 BCE.} Babylonian tablets already contained
              sophisticated arithmetic, quadratic procedures, and examples of
              Pythagorean triples.
        \item \textbf{c. 1650 BCE.} The Rhind Mathematical Papyrus, copied by
              the scribe Ahmes from an older source, recorded Egyptian
              arithmetic and geometric problems.
        \item \textbf{c. 300 BCE.} Euclid organized Greek geometry and number
              theory axiomatically in the thirteen books of
              \textit{Elements}.
        \item \textbf{3rd century BCE.} Archimedes used the method of exhaustion
              to obtain rigorous area and volume results and proved
              \[
                  \frac{223}{71}<\pi<\frac{22}{7}.
              \]
        \item \textbf{263.} Liu Hui completed his influential commentary on
              \textit{The Nine Chapters on the Mathematical Art} and refined
              polygon methods for approximating $\pi$.
        \item \textbf{5th century.} Zu Chongzhi obtained
              $3.1415926<\pi<3.1415927$ and the approximation
              $355/113$.
        \item \textbf{628.} Brahmagupta's
              \textit{Brahmasphutasiddhanta} gave systematic arithmetic rules
              involving zero and negative numbers.
        \item \textbf{c. 825.} Al-Khwarizmi's treatise on
              \emph{al-jabr} gave algebra its name; the Latin form of his name
              produced the word \emph{algorithm}.
        \item \textbf{1202.} Fibonacci's \textit{Liber Abaci} promoted
              Hindu--Arabic numerals and place-value arithmetic in Latin
              Europe.
        \item \textbf{14th--16th centuries.} Madhava and the Kerala school
              developed infinite series for trigonometric functions and $\pi$,
              centuries before the European calculus.
        \item \textbf{1545.} Cardano published \textit{Ars Magna}, containing
              methods for general cubic and quartic equations, based in part on
              work of del Ferro, Tartaglia, and Ferrari.
        \item \textbf{1557.} Robert Recorde introduced the equality sign
              $=$ in \textit{The Whetstone of Witte}.
        \item \textbf{1614.} John Napier published logarithms, transforming
              difficult multiplication and astronomical computation into
              addition.
        \item \textbf{1631.} Thomas Harriot's work, published posthumously,
              introduced the inequality signs $<$ and $>$.
        \item \textbf{1637.} Descartes' \textit{La Géométrie} joined algebra
              and geometry through coordinates and equations.
        \item \textbf{1654.} The Fermat--Pascal correspondence on games of
              chance helped establish mathematical probability.
        \item \textbf{1655.} John Wallis introduced the infinity symbol
              $\infty$ in \textit{De sectionibus conicis}.
        \item \textbf{1665--1666.} During the plague years, Newton developed
              major parts of his calculus of fluxions, series methods,
              mechanics, and gravitation.
        \item \textbf{1675.} Leibniz first used the elongated-$S$ integral sign
              $\int$ in a manuscript; his differential notation
              $\dd x$ and $\dd y$ became standard.
        \item \textbf{1684.} Leibniz published the first paper on differential
              calculus. Newton and Leibniz are now credited with independent
              developments of calculus.
        \item \textbf{1687.} Newton published the
              \textit{Philosophiæ Naturalis Principia Mathematica}, placing
              mechanics and gravitation in a mathematical framework.
        \item \textbf{1706.} William Jones used the symbol $\pi$ for the ratio
              of circumference to diameter; Euler later popularized it.
        \item \textbf{1713 and 1738.} Nicolaus Bernoulli proposed the
              St.\ Petersburg paradox, and Daniel Bernoulli analyzed it using
              expected utility.
        \item \textbf{1715.} Brook Taylor published the general expansion now
              called Taylor's theorem.
        \item \textbf{1730s.} Euler adopted $e$ for the base of natural
              logarithms and $f(x)$ for a function; these notations became
              standard through his work.
        \item \textbf{1736.} Euler solved the Basel problem,
              \[
                  \sum_{n=1}^{\infty}\frac{1}{n^2}=\frac{\pi^2}{6},
              \]
              and his Königsberg bridges paper founded graph theory and helped
              initiate topology.
        \item \textbf{1742.} Christian Goldbach wrote to Euler about the
              conjecture now usually stated as: every even integer greater than
              $2$ is a sum of two primes.
        \item \textbf{1755.} Euler used capital sigma $\Sigma$ for summation in
              \textit{Institutiones calculi differentialis}.
        \item \textbf{1763.} Bayes' essay on inverse probability was published
              posthumously by Richard Price.
        \item \textbf{1777.} Euler used $i$ for $\sqrt{-1}$ in a manuscript
              published in 1794; Gauss later helped make the notation standard.
        \item \textbf{1799.} Gauss's doctoral dissertation supplied the first
              substantially complete proof of the fundamental theorem of
              algebra.
        \item \textbf{1801.} Gauss published
              \textit{Disquisitiones Arithmeticae}, reorganizing number theory;
              in the same year his least-squares calculations helped recover
              the orbit of Ceres.
        \item \textbf{1821.} Cauchy's \textit{Cours d'Analyse} made limits,
              convergence, and continuity central to rigorous analysis.
        \item \textbf{1824.} Abel proved that the general quintic equation
              cannot be solved by radicals.
        \item \textbf{1829--1832.} Lobachevsky and Bolyai published independent
              constructions of hyperbolic geometry, breaking the exclusive
              hold of Euclidean geometry.
        \item \textbf{1830s--1846.} Galois connected polynomial equations with
              permutation groups; Liouville published Galois' manuscripts
              posthumously in 1846.
        \item \textbf{1843.} Hamilton discovered the quaternions, with
              \[
                  i^2=j^2=k^2=ijk=-1.
              \]
        \item \textbf{1844.} Liouville proved that transcendental numbers exist,
              providing the first explicit examples.
        \item \textbf{1847.} Boole published
              \textit{The Mathematical Analysis of Logic}, treating logic by
              algebraic operations.
        \item \textbf{1854.} Riemann's habilitation lecture introduced a broad
              theory of manifolds and curvature, transforming geometry.
        \item \textbf{1858.} August Ferdinand Möbius and Johann Benedict
              Listing independently described the one-sided surface now called
              the Möbius strip.
        \item \textbf{1872.} Dedekind published Dedekind cuts and Cantor
              published a construction of the real numbers, helping complete
              the arithmetization of analysis.
        \item \textbf{1873.} Charles Hermite proved that $e$ is transcendental.
        \item \textbf{1874.} Cantor proved that the real algebraic numbers are
              countable and that transcendental numbers are uncountable.
        \item \textbf{1882.} Ferdinand von Lindemann proved that $\pi$ is
              transcendental. Consequently, squaring the circle with straightedge
              and compass is impossible.
        \item \textbf{1889.} Peano published axioms for the natural numbers in
              a symbolic logical language.
        \item \textbf{1895.} Poincaré published \textit{Analysis Situs}, a
              foundational work of algebraic topology.
        \item \textbf{1896.} Hadamard and de la Vallée Poussin independently
              proved the prime number theorem,
              \[
                  \pi(x)\sim\frac{x}{\ln x}.
              \]
        \item \textbf{1899.} Hilbert's
              \textit{Foundations of Geometry} gave a modern axiomatic treatment
              of Euclidean geometry.
        \item \textbf{1900.} At the second ICM in Paris, Hilbert presented ten
              of the twenty-three problems in his published list; the list
              guided much of twentieth-century mathematics.
        \item \textbf{1901.} Russell discovered the paradox of the set of all
              sets that are not members of themselves, exposing a crisis in
              naive set theory.
        \item \textbf{1904--1908.} Zermelo used the axiom of choice to prove the
              well-ordering theorem and then proposed an axiomatization of set
              theory, later refined into ZF and ZFC.
        \item \textbf{1910--1913.} Whitehead and Russell published the three
              volumes of \textit{Principia Mathematica}.
        \item \textbf{1918.} Emmy Noether proved that continuous symmetries of
              an action correspond to conservation laws.
        \item \textbf{1922.} Banach's work established complete normed spaces
              as central objects of functional analysis; his contraction
              principle became a standard existence theorem.
        \item \textbf{1930.} The Institute for Advanced Study was founded in
              Princeton as an independent research institution without degree
              programs.
        \item \textbf{1931.} Gödel proved his incompleteness theorems: every
              sufficiently strong consistent effectively axiomatized formal
              system is incomplete and cannot prove its own consistency.
        \item \textbf{1933.} Kolmogorov axiomatized probability using measure
              theory.
        \item \textbf{1935.} The Bourbaki group began publishing under the
              fictional name Nicolas Bourbaki; mathematicians of the Lwów
              school began recording problems in \textit{The Scottish Book}.
        \item \textbf{1936.} Turing formalized computation through the Turing
              machine and proved the undecidability of the halting problem.
              The first Fields Medals were awarded in Oslo.
        \item \textbf{1940 and 1963.} Gödel proved that CH is consistent with
              ZF if ZF is consistent; Cohen introduced forcing and proved that
              CH is independent of ZFC, subject to consistency.
        \item \textbf{1947.} George Dantzig introduced the simplex method for
              linear programming.
        \item \textbf{1948.} Shannon founded information theory by defining
              entropy and channel capacity mathematically.
        \item \textbf{1950.} Nash introduced the equilibrium concept now called
              Nash equilibrium.
        \item \textbf{1956.} Milnor discovered exotic smooth structures on the
              $7$-sphere.
        \item \textbf{1958.} The Institut des Hautes Études Scientifiques was
              founded near Paris as a research institute without degree
              programs.
        \item \textbf{1963.} The Atiyah--Singer index theorem connected the
              analytic index of an elliptic operator with topological data.
        \item \textbf{1970.} Matiyasevich completed earlier work of Davis,
              Putnam, and Robinson, proving that Hilbert's tenth problem has
              no algorithmic solution.
        \item \textbf{1976.} Appel and Haken proved the Four Color Theorem with
              extensive computer assistance, creating a landmark debate about
              computer-assisted proof.
        \item \textbf{1982.} The Mathematical Sciences Research Institute was
              founded in Berkeley; it was renamed the Simons Laufer
              Mathematical Sciences Institute in 2022.
        \item \textbf{1983.} Faltings proved the Mordell conjecture: a smooth
              projective curve of genus greater than $1$ over a number field
              has only finitely many rational points.
        \item \textbf{1985.} Louis de Branges proved the Bieberbach conjecture
              on coefficients of univalent functions.
        \item \textbf{1994--1995.} Wiles, with the crucial repair developed
              with Richard Taylor, proved Fermat's Last Theorem through the
              modularity of semistable elliptic curves.
        \item \textbf{1998 and 2014.} Hales announced a computer-assisted proof
              of the Kepler conjecture; the Flyspeck project later completed a
              formal verification.
        \item \textbf{2002--2003.} Perelman posted papers completing Hamilton's
              Ricci-flow program and proving the Poincaré and geometrization
              conjectures.
        \item \textbf{2004.} Green and Tao proved that the prime numbers contain
              arbitrarily long arithmetic progressions.
        \item \textbf{2013.} Yitang Zhang proved the first finite bound on gaps
              between infinitely many pairs of primes; Maynard, Tao, and the
              Polymath project rapidly improved bounded-gap methods.
        \item \textbf{2014.} ICM Seoul was held in South Korea, and Maryam
              Mirzakhani became the first woman to receive the Fields Medal.
        \item \textbf{2016.} Maryna Viazovska proved the optimality of the
              $E_8$ lattice sphere packing in dimension $8$; joint work soon
              settled dimension $24$ using the Leech lattice.
        \item \textbf{2022.} June Huh received the Fields Medal for work joining
              combinatorics, algebraic geometry, and Hodge theory; Maryna
              Viazovska became the second woman Fields Medalist.
    \end{itemize}

    \begin{fact}[Mathematical institutions]
        The distinction tested in Science Quiz is whether an institution grants
        degrees.
        \begin{itemize}
            \item École normale supérieure is a higher-education institution
                  with degree-granting programs.
            \item IAS, IHES, and SLMath, formerly MSRI, are research institutes,
                  not degree-granting universities.
            \item Collège de France offers public lectures and research chairs
                  but does not grant degrees.
        \end{itemize}
    \end{fact}

    \begin{fact}[Königsberg and Kaliningrad]
        Königsberg, home of the Seven Bridges problem and birthplace of Hilbert,
        became Kaliningrad after World War II. Kaliningrad is now part of
        Russia.
    \end{fact}

    \section*{Mathematical Anecdotes}

    Some famous stories are documented, while others survive as traditional
    legends. A story whose details are uncertain is explicitly identified
    below; the uncertainty is part of the fact to remember.

    \begin{itemize}
        \item \textbf{Pythagorean irrationality legend.} A famous but
              historically uncertain legend says that Hippasus was drowned
              after revealing the irrationality of $\sqrt{2}$. The discovery
              itself shattered the belief that every magnitude is a ratio of
              integers.
        \item \textbf{No royal road.} Proclus reports that when Ptolemy asked
              for an easier route to geometry, Euclid replied that there is no
              royal road to geometry.
        \item \textbf{Archimedes and the crown.} Vitruvius tells the traditional
              story that Archimedes recognized how displacement could test a
              crown's purity and ran out crying ``Eureka!'' The precise details
              are not independently verified.
        \item \textbf{Archimedes' circles.} Ancient accounts say that
              Archimedes was killed during the Roman capture of Syracuse while
              absorbed in a geometric diagram. The familiar wording ``Do not
              disturb my circles'' is traditional.
        \item \textbf{Archimedes' tomb.} Cicero reported rediscovering the tomb
              of Archimedes by recognizing the carved sphere and cylinder that
              commemorated Archimedes' favorite result.
        \item \textbf{Cardano and Tartaglia.} Tartaglia disclosed his cubic
              method to Cardano under a promise of secrecy. Cardano later
              published it in \textit{Ars Magna} after learning that del Ferro
              had found the method earlier, beginning a celebrated priority
              dispute.
        \item \textbf{Fermat's margin.} Fermat wrote that he had a marvelous
              proof of his last theorem but that the margin was too small to
              contain it. No valid general proof by Fermat is known.
        \item \textbf{Newton's apple.} Newton later told William Stukeley that
              seeing an apple fall prompted thoughts about gravity. The story
              does not require that an apple struck Newton's head.
        \item \textbf{The calculus priority dispute.} Newton and Leibniz
              developed calculus independently, but a bitter priority dispute
              divided British and continental mathematics for decades.
        \item \textbf{Euler's eyesight.} Euler continued producing mathematics
              at an extraordinary rate after losing almost all of his sight,
              dictating long calculations from memory.
        \item \textbf{Gauss and the arithmetic sum.} A traditional schoolroom
              story says that young Gauss immediately summed
              $1+2+\cdots+100$ by pairing terms to obtain $5050$. The details
              vary among later accounts.
        \item \textbf{Sophie Germain's pseudonym.} Because women were excluded
              from the École Polytechnique, Germain studied its lecture notes
              and submitted work under the name ``M.\ LeBlanc.'' Lagrange
              became a supporter after learning her identity.
        \item \textbf{Galois' last night.} Galois wrote urgent mathematical
              notes before the duel that killed him at age $20$. He did not
              create all of Galois theory in one night; the notes summarized
              and revised ideas developed earlier.
        \item \textbf{Hamilton's bridge.} On October 16, 1843, Hamilton carved
              the quaternion relations
              \[
                  i^2=j^2=k^2=ijk=-1
              \]
              into Brougham Bridge in Dublin after the key idea came to him
              during a walk.
        \item \textbf{Hilbert's problems.} Hilbert's Paris address did not
              orally present all twenty-three problems; ten were delivered in
              the lecture, while the complete list appeared in print.
        \item \textbf{Hilbert's hotel.} The infinite-hotel thought experiment,
              popularized by George Gamow, illustrates that a countably infinite
              full hotel can still accommodate more guests by moving each
              guest to another numbered room.
        \item \textbf{Hilbert and Noether.} When objections were raised to
              Noether teaching at Göttingen because she was a woman, Hilbert
              reportedly replied that the university was not a bathhouse.
        \item \textbf{The taxi-cab number.} Hardy told Ramanujan that taxi
              number $1729$ seemed dull. Ramanujan immediately answered that it
              is the smallest number expressible as a sum of two positive cubes
              in two ways:
              \[
                  1729=1^3+12^3=9^3+10^3.
              \]
        \item \textbf{Nicolas Bourbaki.} Bourbaki is not one mathematician but
              the collective pseudonym of a group that sought a unified,
              axiomatic presentation of modern mathematics.
        \item \textbf{\textit{The Scottish Book}.} Mathematicians of the Lwów
              school wrote problems in a notebook kept at the Scottish Café.
              Prizes ranged from coffee or beer to a live goose.
        \item \textbf{Mazur's goose.} Stanisław Mazur offered a live goose for
              Problem 153 in 1936. Per Enflo solved the basis problem in 1972,
              and Mazur presented the promised goose in a televised ceremony.
        \item \textbf{Gödel's citizenship hearing.} Before his 1947 US
              citizenship examination, Gödel found what he believed was a
              logical route to dictatorship in the Constitution. Einstein and
              Oskar Morgenstern accompanied him and helped prevent the hearing
              from becoming a constitutional debate.
        \item \textbf{Dantzig's homework.} George Dantzig arrived late to Jerzy
              Neyman's statistics class, copied two unsolved problems from the
              board believing they were homework, and later submitted
              solutions to both.
        \item \textbf{Erd\H{o}s' language.} Erd\H{o}s called children
              ``epsilons'' and imagined an ideal book in which God keeps the
              most elegant proofs. The Erd\H{o}s number measures collaboration
              distance from him.
        \item \textbf{Turing's apple.} Turing died from cyanide poisoning in
              1954, and a bitten apple was found nearby. The popular claim that
              the apple was deliberately poisoned, or that it inspired the
              Apple logo, is not established.
        \item \textbf{Milnor's exotic spheres.} Milnor discovered that a
              topological $7$-sphere can carry differentiable structures not
              equivalent to the standard one. The group of oriented homotopy
              $7$-spheres has order $28$.
        \item \textbf{The computer proof controversy.} The 1976 proof of the
              Four Color Theorem required computer checking of many cases,
              forcing mathematicians to debate what it means for a proof to be
              humanly inspectable.
        \item \textbf{Wiles' secret work.} Wiles worked privately for years on
              Fermat's Last Theorem. After announcing a proof in 1993, a gap was
              found; Wiles and Richard Taylor repaired it in 1994.
        \item \textbf{Perelman's refusals.} Perelman declined the 2006 Fields
              Medal and the Clay Millennium Prize for the Poincaré conjecture.
        \item \textbf{Villani's public image.} Cédric Villani became recognizable
              for a cravat and spider brooch. His book
              \textit{Théorème vivant} recounts work with Clément Mouhot on
              nonlinear Landau damping.
        \item \textbf{Birkar's stolen medal.} Caucher Birkar's Fields Medal was
              stolen shortly after the 2018 ceremony in Rio de Janeiro; the IMU
              later presented him with a replacement.
        \item \textbf{No Nobel Prize in Mathematics.} The popular story that
              Nobel excluded mathematics because of a romantic rivalry is
              unsupported. Nobel's will simply named physics, chemistry,
              physiology or medicine, literature, and peace.
    \end{itemize}

    \chapter{Famous Problems and Conjectures}

    Famous problems are usually not tested by asking for complete proofs. They
    are tested through the following information: the standard statement, the
    person or group attached to the problem, whether it is solved, and the main
    idea or keyword attached to the solution. The goal of this chapter is to
    make those associations immediate.

    \section*{Millennium Prize Problems}

    The Clay Mathematics Institute announced the seven Millennium Prize Problems
    in 2000. Each problem carries a prize of one million US dollars. Among the
    seven, only the Poincaré conjecture has been solved.

    \vspace{1em}

    {\small
    \renewcommand{\arraystretch}{1.18}
    \begin{longtable}{@{}L{0.24\textwidth}L{0.30\textwidth}L{0.16\textwidth}L{0.22\textwidth}@{}}
        \toprule
        Problem & Standard mathematical form & Status & Quiz anchors \\
        \midrule
        \endfirsthead

        \toprule
        Problem & Standard mathematical form & Status & Quiz anchors \\
        \midrule
        \endhead

        \midrule
        \multicolumn{4}{r}{Continued on the next page.} \\
        \endfoot

        \bottomrule
        \endlastfoot

        Poincaré Conjecture
        & Every closed simply connected $3$-manifold is homeomorphic to $S^3$.
        & Solved.
        & Poincaré; Perelman; Ricci flow. \\

        Riemann Hypothesis
        & Every nontrivial zero of $\zeta(s)$ has real part $1/2$.
        & Open.
        & Riemann; zeta function; prime numbers. \\

        P versus NP
        & Is $P=NP$?
        & Open.
        & Cook, Karp; computational complexity. \\

        Birch and Swinnerton-Dyer Conjecture
        & For an elliptic curve $E$, $\rk E(\Q)$ equals the order of
          vanishing of $L(E,s)$ at $s=1$.
        & Open.
        & Elliptic curves; $L$-functions. \\

        Hodge Conjecture
        & Rational Hodge classes on smooth projective complex varieties should
          come from algebraic cycles.
        & Open.
        & Hodge theory; algebraic geometry. \\

        Navier--Stokes Existence and Smoothness
        & Smooth initial data for the $3$-dimensional incompressible
          Navier--Stokes equations should have global smooth solutions, or one
          must find a breakdown.
        & Open.
        & Fluid mechanics; PDE; turbulence. \\

        Yang--Mills Existence and Mass Gap
        & Construct quantum Yang--Mills theory on $\R^4$ for compact simple gauge
          groups and prove a positive mass gap.
        & Open.
        & Gauge theory; quantum field theory; mass gap. \\
    \end{longtable}
    }

    \begin{fact}[Millennium Prize status check]
        Among the seven Millennium Prize Problems, the Poincaré conjecture is
        the only one that has been solved. The remaining six are the Riemann
        Hypothesis, P versus NP, the Birch and Swinnerton-Dyer conjecture, the
        Hodge conjecture, Navier--Stokes existence and smoothness, and
        Yang--Mills existence and mass gap.
    \end{fact}

    \begin{theorem}[Poincaré Conjecture]
        Every closed simply connected $3$-manifold is homeomorphic to the
        $3$-sphere $S^3$.
    \end{theorem}

    \begin{remark}
        Henri Poincaré posed the problem in the early twentieth century. Grigori
        Perelman proved it using Ricci flow and ideas from Hamilton's program.
        The key Science Quiz association is
        \[
            \text{Poincaré conjecture}
            \quad\longleftrightarrow\quad
            \text{Perelman}
            \quad\longleftrightarrow\quad
            \text{Ricci flow}.
        \]
    \end{remark}

    \begin{hypothesis}[Riemann Hypothesis]
        Let
        \[
            \zeta(s)=\sum_{n=1}^{\infty}\frac{1}{n^s}
        \]
        be the Riemann zeta function, initially defined for $\Real(s)>1$ and
        continued meromorphically. The Riemann Hypothesis states that every
        nontrivial zero of $\zeta(s)$ has real part $1/2$.
    \end{hypothesis}

    \begin{idea}
        The zeta function is connected to prime numbers by Euler's product
        \[
            \zeta(s)=\prod_{p}\frac{1}{1-p^{-s}}.
        \]
        Thus RH is a statement about the hidden regularity of the distribution
        of prime numbers.
    \end{idea}

    \begin{openproblem}[P versus NP]
        Determine whether
        \[
            P=NP.
        \]
        Here $P$ is the class of decision problems solvable in polynomial time,
        and $NP$ is the class of decision problems whose proposed solutions can
        be verified in polynomial time.
    \end{openproblem}

    \begin{remark}
        The usual slogan is: can every solution that can be checked quickly also
        be found quickly? The keywords are Cook, Karp, NP-completeness, and
        theoretical computer science.
    \end{remark}

    \begin{conjecture}[Birch and Swinnerton-Dyer Conjecture]
        For an elliptic curve $E$ over $\Q$, $\rk E(\Q)$ equals the order
        of vanishing of its $L$-function $L(E,s)$ at $s=1$.
    \end{conjecture}

    \begin{remark}
        The problem connects rational points on elliptic curves with analytic
        behavior of $L$-functions. The phrase ``$\rk$ equals order of vanishing''
        is the main quiz memory.
    \end{remark}

    \begin{conjecture}[Hodge Conjecture]
        On a smooth projective complex variety, every rational cohomology class
        of type $(p,p)$ should be a rational linear combination of classes of
        algebraic cycles.
    \end{conjecture}

    \begin{remark}
        This is one of the deepest bridges between topology and algebraic
        geometry. The short memory is
        \[
            \text{Hodge classes} \quad\longleftrightarrow\quad
            \text{algebraic cycles}.
        \]
    \end{remark}

    \begin{openproblem}[Navier--Stokes Existence and Smoothness]
        For the $3$-dimensional incompressible Navier--Stokes equations,
        determine whether smooth initial data always lead to global smooth
        solutions, or whether finite-time singularities can occur.
    \end{openproblem}

    \begin{remark}
        This is a PDE problem motivated by fluid mechanics. The quiz keywords
        are incompressible fluid, global regularity, singularity, and turbulence.
    \end{remark}

    \begin{openproblem}[Yang--Mills Existence and Mass Gap]
        Construct a mathematically rigorous quantum Yang--Mills theory on
        $\R^4$ for every compact simple gauge group and prove that it has a
        positive mass gap.
    \end{openproblem}

    \begin{remark}
        The problem comes from quantum field theory. In quiz form, the key words
        are gauge theory, quantum field theory, and mass gap.
    \end{remark}

    \section*{Landmark Problems}

    The following problems are not all open problems. Some are solved theorems,
    but they remain famous because of their history, statements, or solution
    methods.

    \vspace{1em}

    {\small
    \renewcommand{\arraystretch}{1.18}
    \begin{longtable}{@{}L{0.23\textwidth}L{0.30\textwidth}L{0.20\textwidth}L{0.19\textwidth}@{}}
        \toprule
        Problem & Standard form & People and status & Quiz anchor \\
        \midrule
        \endfirsthead

        \toprule
        Problem & Standard form & People and status & Quiz anchor \\
        \midrule
        \endhead

        \midrule
        \multicolumn{4}{r}{Continued on the next page.} \\
        \endfoot

        \bottomrule
        \endlastfoot

        Fermat's Last Theorem
        & $x^n+y^n=z^n$ has no positive integer solutions for $n\ge 3$.
        & Fermat; proved by Wiles with Taylor--Wiles.
        & Elliptic curves and modular forms. \\

        Four Color Theorem
        & Every planar map can be colored with at most four colors.
        & Appel--Haken; later Robertson--Sanders--Seymour--Thomas.
        & First famous computer-assisted proof. \\

        Kepler Conjecture
        & The densest packing of congruent spheres in $\R^3$ is the face-centered
          cubic packing.
        & Kepler; proved by Hales.
        & Sphere packing; Flyspeck project. \\

        Thurston Geometrization Conjecture
        & Closed orientable $3$-manifolds decompose into pieces modeled on eight
          geometries.
        & Thurston; completed through Perelman's work.
        & Eight Thurston geometries. \\

        Continuum Hypothesis
        & There is no cardinality strictly between $|\N|$ and $|\R|$.
        & Cantor; independent of ZFC by Gödel and Cohen.
        & Forcing; independence. \\

        Hilbert's Problems
        & A list of $23$ problems proposed in 1900.
        & Hilbert; many solved, some open or independent.
        & Roadmap for twentieth-century mathematics. \\

        Goldbach Conjecture
        & Every even integer greater than $2$ is a sum of two primes.
        & Goldbach; open.
        & Additive number theory. \\

        Twin Prime Conjecture
        & There are infinitely many primes $p$ such that $p+2$ is also prime.
        & Open; bounded gaps by Zhang, Maynard, Tao, Polymath.
        & Prime gaps. \\

        Collatz Conjecture
        & Iterating $n\mapsto n/2$ for even $n$ and $n\mapsto 3n+1$ for odd $n$
          eventually reaches $1$.
        & Collatz; open.
        & Elementary iteration; nonlinear dynamics. \\

        $abc$ Conjecture
        & If $a+b=c$ and $a,b,c$ are coprime, then $c$ is controlled by the
          radical of $abc$, up to $\eps$.
        & Masser--Oesterlé; open in the standard sense.
        & Radical, Diophantine equations. \\
    \end{longtable}
    }

    \begin{theorem}[Fermat's Last Theorem]
        For every integer $n\ge 3$, there do not exist positive integers
        $x,y,z$ such that
        \[
            x^n+y^n=z^n.
        \]
    \end{theorem}

    \begin{remark}
        Fermat wrote that he had a marvelous proof, but no proof appeared in
        his margin. Andrew Wiles proved the theorem through the modularity of
        semistable elliptic curves. The fast association is
        \[
            \text{Fermat's Last Theorem}
            \quad\longleftrightarrow\quad
            \text{Wiles}
            \quad\longleftrightarrow\quad
            \text{elliptic curves and modular forms}.
        \]
    \end{remark}

    \begin{theorem}[Four Color Theorem]
        Every planar map can be colored with at most four colors so that any two
        adjacent regions have different colors.
    \end{theorem}

    \begin{remark}
        The theorem was proved by Kenneth Appel and Wolfgang Haken in 1976 with
        extensive computer assistance. This made it one of the most famous early
        examples of a computer-assisted proof.
    \end{remark}

    \begin{theorem}[Kepler Conjecture]
        The maximum density of a packing of congruent spheres in $\R^3$ is
        attained by the face-centered cubic packing and equals
        \[
            \frac{\pi}{3\sqrt{2}}.
        \]
    \end{theorem}

    \begin{remark}
        Thomas Hales proved the Kepler conjecture using a large computer-assisted
        argument. The Flyspeck project later produced a formal verification of
        the proof. In higher-dimensional sphere packing, Maryna Viazovska proved
        the optimality of the $E_8$ lattice in dimension $8$.
    \end{remark}

    \begin{theorem}[Thurston Geometrization Theorem]
        Every closed orientable $3$-manifold can be decomposed into pieces that
        admit one of the eight Thurston geometries:
        \[
            S^3,\ E^3,\ H^3,\ S^2\times\R,\ H^2\times\R,
            \ \widetilde{\SL}_2(\R),\ \mathrm{Nil},\ \mathrm{Sol}.
        \]
    \end{theorem}

    \begin{remark}
        The Poincaré conjecture is a special case of the geometrization program.
        The 2018 Science Quiz asked for the number of Thurston geometries; the
        answer is $8$.
    \end{remark}

    \begin{hypothesis}[Continuum Hypothesis]
        There is no set whose cardinality is strictly between the cardinality of
        the natural numbers and the cardinality of the real numbers. Equivalently,
        \[
            2^{\aleph_0}=\aleph_1.
        \]
    \end{hypothesis}

    \begin{remark}
        Kurt Gödel proved that CH cannot be disproved from ZFC if ZFC is
        consistent. Paul Cohen proved that CH cannot be proved from ZFC if ZFC
        is consistent. Cohen introduced forcing and received the Fields Medal in
        1966.
    \end{remark}

    \begin{fact}[Hilbert's Problems]
        David Hilbert presented $23$ problems at the 1900 International Congress
        of Mathematicians in Paris. They shaped much of twentieth-century
        mathematics. Useful anchors include:
        \begin{itemize}
            \item Hilbert's first problem $\leftrightarrow$ continuum hypothesis.
            \item Hilbert's second problem $\leftrightarrow$ consistency of arithmetic.
            \item Hilbert's tenth problem $\leftrightarrow$ Diophantine equations;
                  negative solution by Matiyasevich, building on Davis--Putnam--Robinson.
        \end{itemize}
    \end{fact}

    \begin{openproblem}[Goldbach Conjecture]
        Every even integer greater than $2$ is a sum of two prime numbers.
    \end{openproblem}

    \begin{remark}
        Goldbach is one of the easiest famous open problems to state. It belongs
        to additive number theory. The weak Goldbach conjecture, concerning sums
        of three odd primes, was proved by Harald Helfgott.
    \end{remark}

    \begin{openproblem}[Twin Prime Conjecture]
        There are infinitely many primes $p$ such that
        \[
            p+2
        \]
        is also prime.
    \end{openproblem}

    \begin{remark}
        The modern quiz keywords are bounded gaps between primes, Yitang Zhang,
        James Maynard, Terence Tao, and the Polymath project. These results do
        not prove the twin prime conjecture, but they prove that infinitely many
        prime pairs occur within some fixed bounded distance.
    \end{remark}

    \begin{openproblem}[Collatz Conjecture]
        Start with a positive integer $n$. If $n$ is even, replace it by $n/2$.
        If $n$ is odd, replace it by $3n+1$. The Collatz conjecture states that
        repeated iteration always eventually reaches $1$.
    \end{openproblem}

    \begin{remark}
        This problem is famous because the rule is elementary but the global
        behavior is still unknown. It is also called the $3n+1$ problem.
    \end{remark}

    \begin{conjecture}[$abc$ Conjecture]
        For every $\eps>0$, there exists a constant $K_\eps$ such
        that, whenever $a,b,c$ are pairwise coprime positive integers satisfying
        $a+b=c$, we have
        \[
            c<K_\eps \rad(abc)^{1+\eps},
        \]
        where $\rad(n)$ is the product of the distinct prime
        divisors of $n$.
    \end{conjecture}

    \begin{remark}
        The conjecture is powerful because it predicts that the additive
        relation $a+b=c$ strongly restricts the repeated prime factors of
        $abc$. The keyword is ``radical.''
    \end{remark}

    \section*{Combinatorics}

    Combinatorial problems are often elementary to state but resistant to
    direct attack. Their statements involve graphs, set systems, matrices, and
    discrete structures; their solutions frequently import ideas from algebra,
    geometry, probability, or topology.

    \begin{fact}[June Huh and combinatorial log-concavity]
        June Huh's work connects combinatorics with algebraic geometry and Hodge
        theory. The central associations are:
        \begin{itemize}
            \item Read's conjecture $\leftrightarrow$ chromatic polynomials.
            \item Brylawski's and Okounkov's conjectures
                  $\leftrightarrow$ matroid characteristic polynomials and
                  log-concavity.
            \item Rota's conjecture $\leftrightarrow$ log-concavity for matroids.
            \item Mason's conjecture $\leftrightarrow$ independent sets of matroids.
            \item Hodge theory $\leftrightarrow$ the geometric method behind these
                  combinatorial log-concavity results.
        \end{itemize}
        In a quiz, if June Huh appears with a list of conjectures, the odd one
        out is often a famous conjecture from a different field, such as the
        Poincaré conjecture.
    \end{fact}

    \begin{conjecture}[Hadwiger Conjecture]
        Every finite graph $G$ with chromatic number $\chi(G)=k$ contains the
        complete graph $K_k$ as a minor.
    \end{conjecture}

    \begin{remark}
        The conjecture is known for $k\leq 6$ and remains open in general. The
        case $k=5$ is equivalent to the Four Color Theorem.
    \end{remark}

    \begin{conjecture}[Hadamard Conjecture]
        For every positive integer $n$, there exists a $4n\times 4n$ matrix
        $\mathbf{H}$ with entries in $\{-1,1\}$ such that
        \[
            \mathbf{H}\mathbf{H}^{\mathsf T}
            =
            4n\mathbf{I}_{4n}.
        \]
    \end{conjecture}

    \begin{remark}
        A matrix satisfying this orthogonality condition is a Hadamard matrix.
        The necessary divisibility condition is known, but existence for every
        positive multiple of $4$ remains open.
    \end{remark}

    \begin{conjecture}[Union-Closed Sets Conjecture]
        Let $\mathcal{F}$ be a finite nonempty family of finite sets such that
        \[
            A,B\in\mathcal{F}
            \implies
            A\cup B\in\mathcal{F}.
        \]
        Then some element belongs to at least half of the sets in
        $\mathcal{F}$.
    \end{conjecture}

    \begin{remark}
        The conjecture was proposed by Péter Frankl in 1979. It is also called
        Frankl's union-closed sets conjecture and remains open.
    \end{remark}

    \begin{conjecture}[Lonely Runner Conjecture]
        Suppose that $n\geq 2$ runners move at distinct constant speeds around a
        unit circle. For each runner, there is a time at which that runner is at
        circular distance at least
        \[
            \frac{1}{n}
        \]
        from every other runner.
    \end{conjecture}

    \begin{remark}
        After subtracting the chosen runner's speed from every speed, the
        problem becomes one of simultaneous Diophantine approximation. It
        remains open in general.
    \end{remark}

    \section*{Number Theory}

    \begin{openproblem}[Odd Perfect Number Problem]
        Determine whether there exists an odd positive integer $n$ satisfying
        \[
            \sigma(n)=2n,
        \]
        where $\sigma(n)$ is the sum of the positive divisors of $n$.
    \end{openproblem}

    \begin{remark}
        Euclid and Euler characterized the even perfect numbers through Mersenne
        primes. No odd perfect number is known, and nonexistence has not been
        proved.
    \end{remark}

    \begin{conjecture}[Legendre Conjecture]
        For every positive integer $n$, there is a prime $p$ satisfying
        \[
            n^2<p<(n+1)^2.
        \]
    \end{conjecture}

    \begin{remark}
        This is stronger than what follows from the prime number theorem and
        remains open.
    \end{remark}

    \begin{conjecture}[Erd\H{o}s--Straus Conjecture]
        For every integer $n\geq 2$, there exist positive integers $x,y,z$ such
        that
        \[
            \frac{4}{n}
            =
            \frac{1}{x}+\frac{1}{y}+\frac{1}{z}.
        \]
    \end{conjecture}

    \begin{remark}
        The conjecture concerns Egyptian-fraction decompositions and is verified
        computationally for enormous ranges, but no general proof is known.
    \end{remark}

    \begin{conjecture}[Beal Conjecture]
        If positive integers satisfy
        \[
            A^x+B^y=C^z,
            \qquad
            x,y,z>2,
        \]
        then $A$, $B$, and $C$ have a common prime factor.
    \end{conjecture}

    \begin{remark}
        Fermat's Last Theorem is the equal-exponent case with pairwise coprime
        bases. The Beal conjecture remains open.
    \end{remark}

    \begin{conjecture}[Schanuel Conjecture]
        If $z_1,\dots,z_n\in\C$ are linearly independent over $\Q$, then
        \[
            \operatorname{trdeg}_{\Q}
            \Q(z_1,\dots,z_n,e^{z_1},\dots,e^{z_n})
            \geq n.
        \]
    \end{conjecture}

    \begin{remark}
        Schanuel's conjecture is a central organizing conjecture in
        transcendental number theory. It would imply many famous results and
        conjectures about algebraic independence.
    \end{remark}

    \begin{openproblem}[Infinitely Many Mersenne Primes]
        Determine whether there are infinitely many primes of the form
        \[
            2^p-1,
        \]
        where $p$ is prime.
    \end{openproblem}

    \begin{remark}
        Every even perfect number has the Euclid--Euler form
        \[
            2^{p-1}(2^p-1)
        \]
        with $2^p-1$ prime. Thus the infinitude of even perfect numbers is
        equivalent to the infinitude of Mersenne primes.
    \end{remark}

    \begin{openproblem}[Perfect Cuboid Problem]
        Determine whether there exists a rectangular box with integer edge
        lengths $a,b,c$ for which all three face diagonals and the space
        diagonal are integers:
        \[
            a^2+b^2,\quad
            b^2+c^2,\quad
            c^2+a^2,\quad
            a^2+b^2+c^2
        \]
        are all perfect squares.
    \end{openproblem}

    \begin{remark}
        Boxes with integer edges and integer face diagonals are known as Euler
        bricks. Requiring the space diagonal to be integral gives the still-open
        perfect cuboid problem.
    \end{remark}

    \begin{theorem}[Catalan--Mihăilescu Theorem]
        The only solution in integers $x,y,a,b>1$ of
        \[
            x^a-y^b=1
        \]
        is
        \[
            3^2-2^3=1.
        \]
    \end{theorem}

    \begin{remark}
        Catalan posed the conjecture in 1844, and Preda Mihăilescu proved it in
        2002. Thus $8$ and $9$ are the only consecutive positive perfect powers
        greater than $1$.
    \end{remark}

    \section*{Geometry and Topology}

    \begin{openproblem}[Smooth Four-Dimensional Poincaré Conjecture]
        Determine whether every smooth closed $4$-manifold that is homotopy
        equivalent to $S^4$ is diffeomorphic to $S^4$.
    \end{openproblem}

    \begin{remark}
        The topological $4$-dimensional Poincaré conjecture was proved by
        Freedman. The smooth $4$-dimensional case remains open, while the other
        dimensions are settled.
    \end{remark}

    \begin{theorem}[Kakeya Conjecture in Three Dimensions]
        If a set $K\subseteq\R^3$ contains a unit line segment in every
        direction, then
        \[
            \dim_{\mathrm{H}}K=3,
        \]
        and its Minkowski dimension is also $3$.
    \end{theorem}

    \begin{remark}
        Besicovitch showed that Kakeya sets may have Lebesgue measure $0$.
        Hong Wang and Joshua Zahl proved the full dimension conjecture in
        $\R^3$ in 2025. The analogous conjecture remains open in dimensions
        $n\geq 4$.
    \end{remark}

    \begin{openproblem}[Inscribed Square Problem]
        Determine whether every Jordan curve in the plane contains four points
        that are the vertices of a square.
    \end{openproblem}

    \begin{remark}
        Also called the Toeplitz square-peg problem, it is known for many
        regular classes of curves but remains open for arbitrary Jordan curves.
    \end{remark}

    \begin{fact}[Kissing Numbers]
        The kissing number $\tau_n$ is the maximum number of nonoverlapping unit
        spheres that can touch one unit sphere in $\R^n$. Famous exact values
        include
        \[
            \tau_2=6,\qquad
            \tau_3=12,\qquad
            \tau_4=24,\qquad
            \tau_8=240,\qquad
            \tau_{24}=196560.
        \]
        Exact values are unknown in most dimensions.
    \end{fact}

    \begin{openproblem}[Hilbert--Smith Conjecture]
        Every locally compact topological group that acts faithfully and
        continuously on a connected manifold should be a Lie group.
    \end{openproblem}

    \begin{remark}
        The problem asks whether manifolds admit genuinely non-Lie continuous
        symmetry groups. It remains open in full generality.
    \end{remark}

    \section*{Algebra and Logic}

    \begin{conjecture}[Jacobian Conjecture]
        Let
        \[
            F=(F_1,\dots,F_n):\C^n\to\C^n
        \]
        be a polynomial map. If
        \[
            \det\left(
                \frac{\partial F_i}{\partial x_j}
            \right)
        \]
        is a nonzero constant, then $F$ has a polynomial inverse.
    \end{conjecture}

    \begin{remark}
        Keller posed the conjecture in 1939. It is open for every
        $n\geq 2$, even though many reductions are known.
    \end{remark}

    \begin{openproblem}[Inverse Galois Problem]
        Determine whether every finite group $G$ occurs as
        \[
            \Gal(K/\Q)
        \]
        for some finite Galois extension $K$ of $\Q$.
    \end{openproblem}

    \begin{remark}
        The problem is solved for many classes of groups but remains open in
        general.
    \end{remark}

    \begin{openproblem}[Hilbert's Tenth Problem over $\Q$]
        Determine whether there is an algorithm that decides, for every
        polynomial
        \[
            f\in\Z[x_1,\dots,x_n],
        \]
        whether the equation $f=0$ has a solution in $\Q^n$.
    \end{openproblem}

    \begin{remark}
        Matiyasevich's theorem gives a negative answer over $\Z$. The analogous
        decision problem over $\Q$ remains open.
    \end{remark}

    \begin{theorem}[Classification of Finite Simple Groups]
        Every finite simple group is one of the following:
        \begin{itemize}
            \item a cyclic group of prime order;
            \item an alternating group $A_n$ with $n\geq 5$;
            \item a finite simple group of Lie type;
            \item one of the $26$ sporadic simple groups.
        \end{itemize}
    \end{theorem}

    \begin{remark}
        The proof is distributed across thousands of pages by many authors.
        The largest sporadic group is the Monster group.
    \end{remark}

    \begin{theorem}[Restricted Burnside Problem]
        For fixed positive integers $d$ and $n$, there are only finitely many
        finite $d$-generated groups of exponent $n$, up to isomorphism.
    \end{theorem}

    \begin{remark}
        Efim Zelmanov solved the restricted Burnside problem using the theory of
        Lie algebras and received the 1994 Fields Medal. The unrestricted
        Burnside problem has a negative answer in general.
    \end{remark}

    \begin{fact}[Langlands Program]
        The Langlands program predicts deep correspondences between arithmetic
        objects such as Galois representations and analytic objects such as
        automorphic representations. It is not one conjecture but a network of
        conjectures linking number theory, representation theory, and geometry.
    \end{fact}

    \section*{Analysis and Dynamics}

    \begin{openproblem}[Three-Dimensional Euler Singularity Problem]
        For the incompressible Euler equations in $\R^3$, determine whether
        smooth finite-energy initial data can develop a singularity in finite
        time.
    \end{openproblem}

    \begin{remark}
        The Euler equations model an ideal fluid without viscosity. Their global
        regularity problem is distinct from the Navier--Stokes Millennium
        problem.
    \end{remark}

    \begin{conjecture}[Falconer Distance Conjecture]
        If a compact set $E\subseteq\R^n$ satisfies
        \[
            \dim_{\mathrm{H}}E>\frac{n}{2},
        \]
        then its distance set
        \[
            \Delta(E)
            =
            \{\|\mathbf{x}-\mathbf{y}\|:\mathbf{x},\mathbf{y}\in E\}
        \]
        has positive Lebesgue measure.
    \end{conjecture}

    \begin{remark}
        The conjecture connects harmonic analysis, geometric measure theory,
        and discrete distance problems. It remains open in general.
    \end{remark}

    \begin{openproblem}[Hilbert's Sixteenth Problem]
        The second part asks for a uniform upper bound, in terms of the degree
        $n$, on the number and arrangement of limit cycles of a planar
        polynomial vector field of degree $n$.
    \end{openproblem}

    \begin{remark}
        Even the quadratic case is not completely understood. The first part of
        Hilbert's sixteenth problem concerns the topology of real algebraic
        curves.
    \end{remark}

    \begin{theorem}[Bieberbach Conjecture]
        If
        \[
            f(z)=z+\sum_{n=2}^{\infty}a_nz^n
        \]
        is analytic and injective on the unit disk, then
        \[
            |a_n|\leq n
            \qquad
            (n\geq 2).
        \]
    \end{theorem}

    \begin{remark}
        Ludwig Bieberbach conjectured the bound in 1916, and Louis de Branges
        proved it in 1985.
    \end{remark}

    \begin{theorem}[Kadison--Singer Problem]
        Every pure state on the diagonal maximal abelian subalgebra of
        bounded operators on $\ell^2$ has a unique extension to a state on all
        bounded operators.
    \end{theorem}

    \begin{remark}
        Kadison and Singer posed the problem in 1959. Marcus, Spielman, and
        Srivastava solved it in 2013 through an equivalent finite-dimensional
        partition problem and the method of interlacing polynomials.
    \end{remark}
